%=============================================================================
\section{System and Implementation}
\label{sec:system}
%=============================================================================

\subsection{Architecture}

The system is organised as five components arranged around a single mutable
object, the belief vector (Fig.~\ref{fig:arch}). A \emph{catalogue compiler} converts a table of
items and categorical attributes into a boolean question matrix, materialising
once what would otherwise be recomputed at every turn. A \emph{reliability
model} supplies a crossover probability for each attribute. A \emph{query
selector} scores every unasked question against the current belief and the
reliability vector and returns an argmax. An \emph{inference engine} applies
Bayes' rule to the observed answer. A \emph{termination controller} decides
between asking again and committing to a prediction.

The design decision worth stating explicitly is that the reliability model is
consumed in two places rather than one. Systems that model answer noise at all
generally do so in the likelihood, where it prevents a wrong answer from
destroying the posterior. Feeding the same quantity into the selector is what
prevents the wrong answer from being solicited in the first place, and by
Proposition~\ref{prop:uniform} it is only worth doing if the quantity varies
across attributes.

%---------------------------------------------------------------- Fig. arch
\begin{figure}[t]
\centering
\begin{tikzpicture}[node distance=2.6mm]
\node[fdata, text width=25mm] (cat) {Item catalogue\\{\tiny $N$ items $\times$ $M$ attributes}};
\node[fproc, text width=25mm, below=of cat] (comp)
  {Catalogue compiler\\{\tiny binning, pruning}};
\node[fdata, text width=25mm, below=of comp] (X)
  {Answer matrix $X$\\{\tiny $N\times|\Qset|$ boolean}};
\node[frel, text width=20mm, right=7mm of comp] (rel)
  {Reliability\\model $\varepsilon_f$\\{\tiny tiers, or measured}};
\node[fproc, text width=25mm, below=6mm of X] (sel)
  {Query selector\\{\tiny $\arg\max$ RAIG}};
\node[fdata, text width=15mm, left=5mm of sel] (bel) {belief $\bb_t$};
\node[fbox, text width=22mm, below=of sel] (q) {question $q_t$};
\node[fuser, text width=14mm, below=of q] (user) {user};
\node[fproc, text width=25mm, below=of user] (inf)
  {Inference engine\\{\tiny Bayes update}};
\node[fdec, text width=13mm, below=3mm of inf] (term) {stop?};
\node[fbox, text width=22mm, below=3mm of term] (out) {MAP prediction};

\draw[far] (cat) -- (comp);
\draw[far] (comp) -- (X);
\draw[far] (X) -- (sel);
\draw[far] (bel) -- (sel);
\draw[far] (sel) -- (q);
\draw[far] (q) -- (user);
\draw[far] (user) -- node[right, font=\tiny, xshift=0.5mm] {$y_t$} (inf);
\draw[far] (inf) -- (term);
\draw[far] (term) -- node[right, font=\tiny, xshift=0.5mm] {yes} (out);
\draw[far] (term.west) -- ++(-9mm,0) node[above, font=\tiny, pos=0.6] {no}
           |- (bel.south);
\draw[frar] (rel.south) .. controls +(0,-14mm) and +(9mm,4mm) .. (sel.east);
\draw[frar] (rel.east) -- ++(3mm,0) |- (inf.east);
\end{tikzpicture}
\caption{System architecture. The reliability model is consumed by both the
selector and the inference engine (dashed); existing systems typically consume
it in the latter only, which is the omission this work addresses.}
\label{fig:arch}
\end{figure}
%----------------------------------------------------------------------

\subsection{Algorithms}

Algorithms~\ref{alg:select}--\ref{alg:eval} give the procedure. Selection is
one matrix--vector product and a constant number of vectorised transforms.
Every selector is a pure function of belief, question matrix and reliability
vector, with no internal state; sessions are driven entirely by an externally
supplied random stream. This is not incidental tidiness. It is what allows the
same target and the same noise realisation to be replayed across every method
under comparison, which in turn is what licenses paired statistical testing.

\begin{algorithm}[t]
\caption{Reliability-aware query selection}
\label{alg:select}
\begin{algorithmic}[1]
\Require belief $\bb$, answer matrix $X \in \{0,1\}^{N \times |\Qset|}$,
per-question crossover $\hat{\boldsymbol\varepsilon}$, asked mask $U$
\Ensure index of the next question
\State $\mathbf{p} \gets \bb^{\!\top} X$
  \Comment{split mass of all questions, one product}
\State $\tilde{\mathbf{p}} \gets \mathbf{p} \odot (1-\hat{\boldsymbol\varepsilon})
   + (1-\mathbf{p}) \odot \hat{\boldsymbol\varepsilon}$
\State $\mathbf{s} \gets \Hb(\tilde{\mathbf{p}}) - \Hb(\hat{\boldsymbol\varepsilon})$
\If{modelling abstention}
  \State $\mathbf{s} \gets (1-\hat{\boldsymbol\delta}) \odot \mathbf{s}$
    \Comment{Eq.~\eqref{eq:erasure}}
\EndIf
\State $s_j \gets -\infty$ for all $j$ with $U_j = 1$
\State \Return $\arg\max_j s_j$
\end{algorithmic}
\end{algorithm}

\begin{algorithm}[t]
\caption{Belief update under BSC$(\hat\varepsilon)$}
\label{alg:update}
\begin{algorithmic}[1]
\Require belief $\bb$, answer column $\bx_q$, observation $y$, crossover $\hat\varepsilon$
\If{$y$ is an erasure} \State \Return $\bb$ \Comment{no information} \EndIf
\State $\ell_c \gets 1-\hat\varepsilon$ if $x_q(c)=y$ else $\hat\varepsilon$, for all $c$
\State $\bb' \gets \bb \odot \boldsymbol\ell$; \quad $Z \gets \sum_c b'_c$
\If{$Z = 0$} \Comment{reachable only at $\hat\varepsilon=0$}
  \State \Return $\bb$ \Comment{stated convention; see Prop.~\ref{prop:limit}}
\EndIf
\State \Return $\bb' / Z$
\end{algorithmic}
\end{algorithm}

\begin{algorithm}[t]
\caption{Identification session}
\label{alg:session}
\begin{algorithmic}[1]
\Require catalogue, target $c^\ast$, budget $T$, selector $\pi$,
assumed $\hat{\boldsymbol\varepsilon}$, true $\boldsymbol\varepsilon$,
random stream $u_{1:T}$
\State $\bb \gets$ uniform; $U \gets \mathbf{0}$
\For{$t = 1$ \textbf{to} $T$}
  \State $q \gets \pi(\bb, X, \hat{\boldsymbol\varepsilon}, U)$; $U_q \gets 1$
  \State $y \gets x_q(c^\ast)$
  \If{$u_t < \varepsilon_{f(q)}$} $y \gets 1-y$ \Comment{answer error} \EndIf
  \State $\bb \gets$ \textsc{Update}$(\bb, \bx_q, y, \hat\varepsilon_{f(q)})$
\EndFor
\State \Return $\arg\max_c \bb(c)$
\end{algorithmic}
\end{algorithm}

\begin{algorithm}[t]
\caption{Paired evaluation with common random numbers}
\label{alg:eval}
\begin{algorithmic}[1]
\Require methods $\mathcal{M}$, trials $R$, seeds $S$, budget $T$
\For{each seed $s \in S$}
  \State draw targets $c^\ast_{1:R}$ and streams $u_{1:R,1:T}$ from seed $s$
  \For{$i = 1$ \textbf{to} $R$}
    \For{each $m \in \mathcal{M}$}
      \State record outcome of \textsc{Session}$(c^\ast_i, u_{i,\cdot}, m)$
    \EndFor
  \EndFor
\EndFor
\State \Return per-trial outcomes for paired inference
\end{algorithmic}
\end{algorithm}

One implementation detail is omitted from Algorithm~\ref{alg:select} because it
never fires in our experiments and stating it inline would obscure the
selection rule. A method restricted to a subset of attributes could in
principle exhaust its allowed pool before the budget is spent; the
implementation then falls back to the unrestricted unasked pool rather than
returning nothing. With $T=20$ and at least $126$ allowed questions for every
method including the restricted baseline, this branch is unreachable in every
run reported here, and the restricted baseline's poor showing is therefore not
an artefact of it.

Sharing $u_{i,\cdot}$ across methods in Algorithm~\ref{alg:eval} is deliberate.
Methods diverge in which questions they ask and therefore in which crossover
probabilities they face, but they consult the same underlying draws, so
differences in outcome are attributable to question choice rather than to
sampling luck. This substantially reduces variance relative to independent
sampling and permits exact paired testing~\cite{mcnemar1947}.

\subsection{Complexity}

Let $|\Qset|$ be the number of compiled questions. Selection is one
matrix--vector product and a constant number of vectorised transforms, hence
$O(N|\Qset|)$ time and $O(|\Qset|)$ working memory; the update is $O(N)$. A
session of $T$ turns costs $O(TN|\Qset|)$, and the compiled matrix occupies
$N|\Qset|$ bits. Table~\ref{tab:complexity} summarises. Critically, RAIG and
classical IG have identical asymptotic complexity, and the constant is
negligible too: the correction adds one elementwise affine map and one
subtraction on a vector of length $|\Qset|$ to a matrix--vector product costing
$O(N|\Qset|)$, so the expected relative overhead is $1/N$, here
$1.3\times10^{-5}$. Section~\ref{sec:results-efficiency} reports what happens
when one tries to measure a difference that small on a shared machine.

\begin{table}[t]
\caption{Complexity per operation. RAIG and classical IG differ only in the
score transform, which is $O(|\Qset|)$ for both.}
\label{tab:complexity}
\centering
\footnotesize
\begin{tabular}{@{}lll@{}}
\toprule
\textbf{Operation} & \textbf{Time} & \textbf{Memory} \\
\midrule
Catalogue compilation & $O(NM\bar{V})$ & $O(N|\Qset|)$ bits \\
Split-mass computation & $O(N|\Qset|)$ & $O(|\Qset|)$ \\
Score transform (IG or RAIG) & $O(|\Qset|)$ & $O(|\Qset|)$ \\
Belief update & $O(N)$ & $O(N)$ \\
Full session & $O(TN|\Qset|)$ & $O(N|\Qset|)$ bits \\
\bottomrule
\end{tabular}
\end{table}

\subsection{Implementation and Reproducibility}

The reference implementation is Python~3.13 against NumPy, pandas and SciPy;
no deep-learning framework is required, and the entire selection and inference
path is dense linear algebra. Catalogue compilation materialises one boolean
column per non-degenerate attribute--value pair; questions true of every item
or of none are discarded at compile time, since they carry zero information
under either criterion.

The engineering decision with the largest downstream effect is the separation
of the assumed crossover $\hat{\boldsymbol\varepsilon}$ from the true crossover
$\boldsymbol\varepsilon$ in the session signature. The simulator draws errors
from the latter; the method sees only the former. Conflating them---the natural
shortcut---would make every reported result circular, because the method would
be evaluated against a user whose behaviour it defines. Keeping them distinct
is what makes the misspecification studies of
Sections~\ref{sec:results-ensemble} and~\ref{sec:results-corruption} possible
at all.

\emph{Determinism.} The BLAS thread count is pinned to one before NumPy is
imported. This is a correctness requirement, not a performance tweak. A
multi-threaded matmul sums $(|\Qset|,N)\times(N,\cdot)$ in a shape-dependent
order, so the split-mass vector differs in the last float32 unit in the last
place between runs with different batch shapes; selection is an argmax over
that vector, near-ties are common in a catalogue with many equal-mass
questions, and a single different question at turn one sends the entire session
down a different path. Pinning to one thread makes every run bitwise
reproducible and---verified empirically as a gate test---makes the per-method
results of a run with $M$ methods identical to those of a run with any subset
of them. That property is what allows a new method to be added to an existing
comparison without re-running the incumbents, and it is worth the modest cost
of stating. On our hardware it is also faster, the matmul being memory-bound.
